In the context of information security, a \textbf{security property} is a characteristic or feature of an application or of one of its components (e.g., a cryptographic algorithm or protocol) that contributes to the protection of data in one or more phases of its lifecycle --- i.e., at rest, in transit, or in use. Security properties may be more or less relevant (or even undesirable) based on the requirements of the underlying scenario. Furthermore, security properties are sometimes opposed (e.g., non-repudiation vs. deniability) and thus cannot be achieved simultaneously. Below, we list security properties common in cryptography.

\paragraph*{Confidentiality.} Data can only be accessed by authorized agents. In cryptography, confidentiality means that a given ciphertext can be decrypted only by the intended recipient(s).

\readBlock{On the difference between confidentiality and privacy}{Confidentiality is not to be confused with the (related yet different) concept of privacy: while confidentiality concerns who can access what data, privacy concerns how data are used to derived information (see \url{https://resources.data.gov/glossary/data-vs.-information/} for a clarification of data vs. information). Traditional approaches to privacy are either \gls{pets} --- usually involving advanced cryptographic constructions (e.g., see \Cref{sec:advanced_cryptography:functional_encryption,sec:advanced_cryptography:homomorphic_encryption}) or confidential computing (\Cref{sec:cryptographic_modules:tee}) --- and data anonymization techniques (not covered in this handbook). In both approaches, there is a trade-off between the degree of privacy achieved within a system and the level of data utility. Moreover, such approaches are (sometimes on purpose) often not sufficient to prevent arbitrary interference. In fact, there may be misuses and ``function creeps'', that is, request of unnecessary data or use of data beyond the intended purpose. Worryingly, \gls{pets} are nowadays being implemented by few big companies, that consequently can decide how to design, use, and update them. Summarizing, \textbf{privacy is not a goal but a mean with which to protect individuals \cite{stadler_why_2022}, and it essentially amounts to data purpose limitation}.}

\paragraph*{Integrity.} Any alteration of some data is easily noticeable --- at least to the intended recipient(s). In general, integrity states that data can be modified by authorized agents only.

\warningBlock{Confusion on the meaning of integrity}{Integrity does not mean that data cannot be tampered with by an attacker. In other words, data can be modified by unauthorized agents even in the presence of cryptographic-based integrity mechanisms (e.g., digital signatures).}

\readBlock{Cybersecurity Integrity vs. Information Integrity}{The meaning of integrity may assume different nuances depending on the underlying context: integrity in cybersecurity mainly concerns digital technology (e.g., hardware and software), while information integrity mainly concerns the users of digital technology (e.g., disinformation, misinformation, and trust). The World Economic Forum's post ``Why cybersecurity and information integrity are two sides of the same coin'' --- available at the link \url{https://www.weforum.org/stories/2025/10/cybersecurity-information-integrity/} --- further discusses this duality.}

\paragraph*{Availability.} Data (and services) are readily available to authorized agents.

\paragraph*{Authenticity.} A specific instance of integrity stating that claims about the origin of data can be verified cryptographically. We can identify (at least) two kinds of authenticity:
\begin{itemize}
	\item \textbf{sender authenticity} (also called \textit{entity authenticity}): a specific instance of authenticity stating that the intended recipient(s) know who created a given message. This is often the case when using asymmetric cryptography (\Cref{sec:asymmetric_cryptography}) to provide authenticity;
	\item \textbf{message authenticity} (also called \textit{data origin authenticity} or \textit{ciphertext authenticity} according to the context): a specific instance of authenticity stating that the intended recipient(s) know that whoever created a given message was authorized (e.g., had access to certain cryptographic material). This is often the case when using symmetric cryptography (\Cref{sec:symmetric_cryptography}) to provide authenticity. Message authenticity is sometimes also called \textit{plausible deniability}.
\end{itemize}

\warningBlock{Authenticity vs. authentication vs. authorization}{Authenticity is not to be confused with the (related yet different) concepts of authentication and authorization: authentication consists in verifying the identity of a party --- more generically, an agent --- while authorization consists in verifying what actions an agent can perform on what resources and accordingly enforcing the resulting decision.}

\paragraph*{Public Verifiability.} Anyone within a system can verify authenticity claims made on data. Specifically, public verifiability requires that no secret data (e.g., private or secret keys) is needed for such verification.

\paragraph*{Non-Repudiation.} a specific instance of public verifiability where the proof of authenticity may only have been created by using some secret data (e.g., private keys) reliably associated with a specific party who, therefore, cannot deny having produced it. Non-repudiation is a stronger security property than sender authenticity. In fact, non-repudiation not only requires that some data can be attributed to the sender (i.e., sender authenticity) but also ensures that this attribution holds up even if the sender disputes it. Hence, non-repudiation implies sender authenticity, but sender authenticity does not imply non-repudiation.
\remarkBlock{Legal aspects of non-repudiation}{In legal terms, non-repudiation typically implies that it is possible to hold one party to their commitments in a legal court. However, technical solutions said to offer non-repudiation often do not hold up in court because of various real-world gaps. For instance:
\begin{itemize}
	\item it may not be provable that the secret data (i.e., key) is correctly associated with the party and that no other party has ever had access to it;
	\item such technical solutions do not adhere to relevant legal frameworks and policies establishing the legitimacy of signatures as evidence;
	\item a system holding the key, if compromised (e.g., a laptop infected by malware), may perform some actions that are not warranted by the associated party.
\end{itemize}		
Usually, a quite complex system of technical and legal assurances is required to complement the cryptographic mechanism for non-repudiation (e.g., digital signatures) and fill the aforementioned gaps.}

\paragraph*{Unforgeability.} An attacker cannot create a ciphertext or a digital signature that an honest verifier deems authentic. In the context of digital signatures (\Cref{sec:digital_signatures}), unforgeability implies that an attacker cannot produce a digital signature that authenticates to --- that is, seems to have been created by --- another party without having the party's private key.
\remarkBlock{Unforgeability vs. non-repudiation}{Unforgeability is a technical security property of cryptographic algorithms, stating that a ciphertext or digital signature that authenticates to a specific symmetric or public key can have been created only by using the same symmetric key or the corresponding private key. Still, unforgeability provides no information on which party used the private or symmetric key --- that is, unforgeability yields no notion of binding between keys and parties. Conversely, non-repudiation is a legal and practical security property related to how ciphertexts and digital signatures are created and verified. Non-repudiation implies unforgeability and also that a ciphertext or digital signature was created by the owner party of the key in a non-repudiable way, where ownership is usually specified through a \gls{pki} --- see \Cref{sec:binding_keys_to_parties}.}

\paragraph*{Deniability.} Whoever created a ciphertext or digital signature can deny it. Deniability may be immediate or may be provided after a certain amount of time or steps in the underlying application (e.g., disclosure of previously secret data like private keys). Deniability is the complement to non-repudiation. In other words, if an application offers deniability, then it cannot offer non-repudiation at the same time --- and vice versa.

\paragraph*{\gls{zk}.} A party (\textbf{prover}) can prove to another party (\textbf{verifier}) that a certain claim (mathematical \textbf{statement}) is true in such a way that the verifier learns nothing beyond the fact that the claim is true. \gls{zk} is a fundamental property of \glspl{zkp} and is discussed in more detail in \Cref{sec:advanced_cryptography:zero_knowledge_proofs}.

\paragraph*{Verifiability.} A party can verify that a cryptographic computation --- the execution of which has been outsourced to a third party (such as a cloud provider) --- has been executed correctly, without the first party having to perform the cryptographic computation themselves. Usually, the first party verifies the correctness of the output of the cryptographic computation non-interactively using a proof of correctness --- often generated with \glspl{zkp} or \gls{he} (\Cref{sec:advanced_cryptography:homomorphic_encryption}) --- returned by the third party along with the aforementioned output.

\paragraph*{Forward Secrecy.} The compromise of a long-term asymmetric private key or session key --- where the compromise happens in the future --- does not compromise current or past session keys (and, consequently, the confidentiality of ciphertexts created in the past). Forward secrecy is relevant to data in transit but not to data at rest, where parties need to access (encrypted) stored data at any moment in the future, i.e., relying on a long-term key. Forward secrecy can be split into \textit{sender forward secrecy} --- that is, forward secrecy in case a key is compromised sender-side --- and \textit{recipient forward secrecy} --- that is, forward secrecy in case a key is compromised recipient-side. A commonly adopted approach to achieve forward secrecy is employing single-use ephemeral key pairs in \gls{dhe} (\Cref{sec:diffie_hellman}) --- while long-term key pairs should only be used when (sender) authenticity is required.

\remarkBlock{Perfect forward secrecy}{Perfect secrecy is sometimes also called \textit{perfect forward secrecy}, although some deem the word ``perfect'' to be a misuse in this case.\footnote{\url{https://blog.cryptographyengineering.com/2014/08/13/whats-matter-with-pgp/}}}

\readBlock{Relevance of forward secrecy}{Michael Horowitz's post ``Perfect Forward Secrecy can block the NSA from secure web pages, but no one uses it'' --- available at the link \url{https://www.computerworld.com/article/2473792/perfect-forward-secrecy-can-block-the-nsa-from-secure-web-pages--but-no-one-uses-it.html} --- and Parker Higgins's post ``Pushing for Perfect Forward Secrecy, an Important Web Privacy Protection'' --- available at the link \url{https://www.eff.org/deeplinks/2013/08/pushing-perfect-forward-secrecy-important-web-privacy-protection} --- further discuss the relevance of forward secrecy to avoid mass surveillance.}

\paragraph*{Backward Secrecy.} The compromise of a long-term asymmetric private key or session key --- where the compromise happened in the past --- does not compromise current or future session keys (and, consequently, the confidentiality of ciphertexts created in the future). In other words, backward secrecy implies that key establishment protocols can \textit{self-heal} after the violation of a long-term asymmetric private key or session key. Backward secrecy is also called \textit{future secrecy}.

\remarkBlock{Security properties and \glspl{kga}}{Although \glspl{kga} are typically trusted to manage keys on behalf of other parties, the presence of \glspl{kga} --- according to how \glspl{kga} are actually used --- may invalidate some security properties like authenticity and non-repudiation.}
